%------------------------------------------------------------------------------%
\begin{question}
Encontre o plano tangente ao gráfico da função
\[
  f(x, y) = x^2 + y^2 -2xy -x +3y + 4
\]
no ponto \((2, -3)\)
\end{question}

%------------------------------------------------------------------------------%
\begin{resolution}{Solução}
  A equação do plano tangente ao gráfico da função $f$ no ponto \((2, -3)\) é
  \pause
  \[
    f_x(a, b)\big(x-a\big) + f_y(a, b)\big(y-b\big) - \big(z-f(x, y)\big) = 0
  \]
  \[
    \pause
    f_x(2, -3)\big(x-2\big) + f_y(2, -3)\big(y+3\big) - \big(z-f(2, -3)\big) = 0
  \]
\end{resolution}

%------------------------------------------------------------------------------%
\begin{resolution}{Calculando a derivada parcial em $x$}
  \[
    f_x(x, y)
    \pause
    \;=\; \D{f}{x}(x, y)
    \pause
    \;=\; \D{}{x}\left(x^2 + y^2 -2xy -x +3y + 4\right)
    \pause
    \;=\; 2x -2y -1
  \]

  \pause
  \bigskip
  \[
    f_x(2, -3)
    \pause
    \;=\; \left(2x -2y -1\right)\evalat{(2, -3)}{}
    \pause
    =\; 2(2) - 2(-3) - 1
    \pause
    \;=\; 4 + 6 - 1
    \pause
    \;=\; 9
  \]
\end{resolution}

%------------------------------------------------------------------------------%
\begin{resolution}{Calculando a derivada parcial em $y$}
  \[
    f_y(x, y)
    \pause
    \;=\; \D{f}{y}(x, y)
    \pause
    \;=\; \D{}{y}\left(x^2 + y^2 -2xy -x +3y + 4\right)
    \pause
    \;=\; 2y - 2x + 3
  \]

  \pause
  \bigskip
  \[
    f_y(2, -3)
    \pause
    \;=\; \left(2y - 2x + 3\right)\evalat{(2, -3)}{}
    \pause
    \;=\; 2(-3) - 2(2) + 3
    \pause
    \;=\; -6 -4 + 3
    \pause
    \;=\; -7
  \]
\end{resolution}

%------------------------------------------------------------------------------%
\begin{resolution}{Avaliando $f$ no ponto \((2, -3)\)}
  \begin{align*}
    \onslide<+->{f(2, -3)}
    \onslide<+->{& = \left(x^2 + y^2 -2xy -x +3y + 4\right)\evalat{(2, -3)}{} \\}
    \onslide<+->{& = 2^2 + (-3)^2 -2(2)(-3) -2 +3(-3) + 4 \\[2mm]}
    \onslide<+->{& = 4 + 9 + 12 - 2 - 9 + 4 \\[2mm]}
    \onslide<+->{& = 18}
  \end{align*}
\end{resolution}

%------------------------------------------------------------------------------%
\begin{resolution}{Equação do Plano Tangente}
  \begin{gather*}
      \onslide<+->{f_x(2, -3)\big(x-2\big) + f_y(2, -3)\big(y+3\big) - \big(z-f(2, -3)\big) = 0 \\[3mm]}
      \onslide<+->{9(x-2) -7(y+3) - (z-18) = 0 \\[3mm]}
      \onslide<+->{9x - 18 -7y -21 -z +18 = 0  \\[3mm]}
      \onslide<+->{9x -7y -z = 21}
  \end{gather*}
\end{resolution}

%------------------------------------------------------------------------------%
